\documentclass[11pt]{article}

\usepackage[margin=1in]{geometry}
\usepackage{amsmath}
\usepackage{booktabs}
\usepackage{caption}
\usepackage[skip=6pt, indent=0pt]{parskip}
\usepackage{microtype}
\usepackage[T1]{fontenc}
\usepackage{lmodern}
\usepackage{xcolor}

\newcommand{\clue}[2]{\textsc{#1}~#2}

\title{Codenames Spymaster: Midterm Progress Report}
\author{Shane Lee}
\date{October 2026}

\begin{document}
\maketitle

\section{Setup}

The goal is still the one from the proposal: a model that gives Codenames clues
meant for human guessers, with no language model producing the clue. The
spymaster sees the board, its team's words, the opponent's, the neutrals and the
assassin, and outputs one clue word and a number.

\paragraph{Evaluating models.} I can't play hundreds of games myself for every
model, so the guesser in every game is a language model. It is only ever the
guesser: it never sees the roles and never writes a clue. Each turn it is shown
the clue, the number and the unrevealed words, with this prompt:

\begin{quote}\small
You are playing the guesser role in the board game Codenames. Your spymaster gave
the clue ``\emph{clue}'' for \emph{N} word(s). Here are the words still available to
guess: \emph{[words]}. Rank ALL of the words above by how likely you would guess
them for this clue, most likely first. Respond with ONLY a JSON array of the words
[\ldots]
\end{quote}

The turn then reveals the guesser's top $N$ words in order and stops at the first
word that isn't ours. Under this prompt the guesser always takes exactly $N$
guesses, unless it hits a wrong word first.

At first the guesser was Claude Sonnet 5. That cost several dollars per 100
games, which limited how many comparisons I could run. I switched to
gpt-oss-120b (an open-weight model, served by DeepInfra), with the same prompt.
It costs about \$0.07 per thousand calls, so a full comparison costs
about ten cents.

\paragraph{The arena.} Two models play each other on a fixed set of 100 boards.
Each board is played twice with the sides swapped, so neither model benefits
from going first or from an easy board. The boards are drawn from 150 words that
no model was trained on. Most comparisons are against a fixed opponent, the
current \emph{incumbent}. The headline statistic is a sign test on the boards one
model won in both seatings; the win rate is reported beside it.

\paragraph{Two numbers per model.} Most of my models are split into two parts: a
\emph{listener}, which predicts what a guesser will pick for a given clue, and a
\emph{decision rule}, which searches about 11,000 candidate clues with the
listener and picks the best one. The listener can be scored on its own: I use
McFadden $R^2$ (1 minus the model's log-loss on the guess actually made, over the
log-loss of guessing uniformly among the words left). This is measured on a
held-out set of 4,378 guesses by Claude Sonnet 5.5 on words no model trained on.
Win rate measures the whole spymaster; $R^2$ measures how well it understands the
guesser. The two don't always move together.

\section{The Gaussian model}

The first real model scores a clue by expected reward under an assumed guesser.
It takes the similarity between the clue and each board word in an embedding
space (ConceptNet Numberbatch), standardized per word. It then assumes the
guesser perceives each similarity with independent Gaussian noise of standard
deviation $\sigma$ and picks words in order of the noisy similarity. Under that
assumption the expected reward of announcing $k$ (own words count $+1$; a neutral
costs $0.2$, an opponent word $1$, the assassin $10$) has a closed form, so every
clue and number in the pool can be scored exactly.

$\sigma$ is the only parameter: how much the guesser's sense of similarity
differs from the embedding's. I swept it, against a simple centroid baseline,
with Sonnet guessing, 50 boards each played both ways:

\begin{center}
\begin{tabular}{rrrrr}
\toprule
$\sigma$ & win\% vs centroid & mean number & losses in the race & losses to the assassin \\
\midrule
1.0  & 69 & 2.67 & 17 & 14 \\
1.25 & 70 & 2.27 & 22 & 8 \\
1.5  & 70 & 1.93 & 24 & 6 \\
2.0  & 57 & 1.42 & 41 & 2 \\
2.5  & 51 & 1.16 & 46 & 3 \\
\bottomrule
\end{tabular}
\end{center}

Win rate plateaus between 1.0 and 1.5 and drops sharply above that. Inside the
plateau, the failure mode shifts: a smaller $\sigma$ gives bigger numbers, which
wins more races but loses more games to the assassin. I kept $\sigma=1.5$, which
has the same win rate as 1.0 with half the assassin losses.

The interesting part is that I also fit $\sigma$ directly to how Sonnet actually
ranks words, and got about $2.06$. That is in the bad region. So the $\sigma$ that
\emph{describes} the guesser best is not the one that \emph{plays} best. The
per-turn objective undervalues tempo, and a spymaster that believes the guesser
is a bit sharper than it is gives bigger clues and wins more races. This came up
again later: fitting the guesser better and winning more are related but not the
same thing.

\section{A learned listener (GBT)}

The Gaussian model's guesser is a single noise parameter on one embedding. The
next step was to learn the guesser instead. I bought gpt-oss's full rankings for
about 25,000 positions (47,223 individual picks), many of them on deliberately
bad clues, so the model also learns what a guesser does when a clue doesn't work.
A LightGBM model scores each (clue, board word) pair, and the probability of the
next pick is a softmax over the scores of the words still on the board. It is
trained on that objective directly.

The final version has 44 features for each (clue, board word) pair:
\begin{itemize}
  \item similarities in several embedding spaces;
  \item human free-association norms in both directions (SWOW: does the clue
        make people think of the word, and the reverse);
  \item language-model pointwise mutual information;
  \item concreteness norms and WordNet relations.
\end{itemize}
The association and PMI features carry most of the gain over embeddings alone.

Treating the listener's scores as a Plackett--Luce model, the expected reward of
a clue still has an exact closed form. The search keeps the Gaussian model's top
200 clues and rescores them with the listener.

\textbf{Results.}
\begin{itemize}
  \item As a model of the guesser: held-out $R^2$ 0.346 against the Gaussian
        model's 0.226 (on gpt-oss's rankings of held-out boards).
  \item In games, with Sonnet guessing: 72\% against the Gaussian model at
        $\sigma=1.5$ over 100 games (95\% CI 63--80\%).
  \item Against Claude Sonnet 5 \emph{as the spymaster}, with Sonnet also
        guessing for both sides: 58\% over 200 games (28 boards won both ways
        against 12, $p = 0.017$). Most of the margin came from the assassin:
        Sonnet's clues led its own guesser to the assassin 31 times, against 13.
\end{itemize}

This model (\texttt{learned\_listener}) became the incumbent and the fixed
opponent for everything that follows. On the current Sonnet 5.5 held-out set
its $R^2$ is 0.3205, the baseline for the table in Section~\ref{sec:table}.

\section{Association features from gpt-oss}

Embeddings capture how words are used, but what a guesser picks depends on what
the clue \emph{reminds} them of. So I asked gpt-oss directly, with no board: ``List
the 25 words that the word \emph{clue} most strongly makes you think of''. I asked
five times per word at temperature 1, for every clue in the pool and every board
word. From those lists come two features per (clue, board word) pair:
\begin{itemize}
  \item the share of lists that name the board word;
  \item a reciprocal-rank score, so a word named first counts far more than one
        named 25th.
\end{itemize}
For example, every list for \clue{river}{} names \emph{water} first, and
\emph{bank} and \emph{mouth} appear in their river senses. A second version (the
\emph{association profile}) adds features describing the clue itself:
\begin{itemize}
  \item how vague its associations are (how much the five lists agree);
  \item its rarity and number of senses;
  \item reverse associations: whether each board word's own lists name the clue.
\end{itemize}

These lists are free at play time, since every clue in the pool already has
them. There is an obvious risk in features built from the same model that guesses
in the arena. That is why $R^2$ is measured on Sonnet's guesses, which never saw
these lists.

\textbf{Results.}
\begin{itemize}
  \item $R^2$ rises from 0.3240 to 0.3293 with the associations, and to 0.3309
        with the profile. That is the largest single feature gain on Sonnet.
  \item In games, the association listener alone is not significant (53.1\%,
        $p = 0.43$). The profile version wins 56.7\% against the incumbent
        (22 boards won both ways against 9, $p = 0.029$), and it became the new
        incumbent.
\end{itemize}

\section{The decision rule: winning, not words}

Up to here the decision rule maximizes expected net words for the turn. That
can't see tempo: two words matter more when they leave you one word from
winning. So I estimated $V(a, b)$, the probability that the side to move wins
with $a$ words left against $b$. It comes from 10,266 finished gpt-oss games,
smoothed and forced to be monotone. The rule then picks the clue with the highest
probability of winning after the turn. For each way the turn can end, it takes
the listener's probability of that ending times the $V$ of the position it leaves.

\begin{itemize}
  \item With the incumbent's own listener, the new rule is level with the old
        one (51.5\%).
  \item With the association listener it is the best result on the suite:
        59.1\% (28 boards against 10, $p = 0.005$).
\end{itemize}
The rule uses the listener's probabilities directly, so a better listener matters
more under it. The price is roughly twice the incumbent's assassin losses. This
model gives bigger numbers (mean 2.6 against 2.3), and its bets sometimes fail
badly.

\section{Later picks: temperatures and pick order}

The listener gives each word one score per clue, and every pick of the turn uses
those same scores. But gpt-oss is far less predictable on its later picks. The
incumbent's $R^2$ falls from 0.55 on the first pick to 0.06 on the fourth.
Treating pick 3 with the same confidence as pick 1 makes the spymaster overrate
clues for 3 or 4.

\paragraph{Per-pick temperatures.} Divide the scores at pick $j$ by a fitted
temperature $\tau_j$, with $\tau_1 = 1$. The fitted values grow with the pick (1.09,
1.26, 1.41) and are nearly the same for every listener, so the flattening is a
property of the guesser, not of one model. $R^2$ improves (0.3205 to 0.3271 on the
incumbent's listener, 0.3309 to 0.3372 on the profile listener), all of it in the
later picks. In games, though, the temperature version of the incumbent
\emph{lost}: 44.3\% (8 boards against 19, $p = 0.05$). A better fit made it more
cautious, and caution lost races, which matches what the $\sigma$ sweep showed.

\paragraph{Pick order.} Instead of a fixed flattening, let the model learn what
changes from pick to pick. Two features are added, the pick index $j$ and $j - k$
(how far past the announced number this pick is). Each pick is still a softmax
over the words left. This beats the temperatures slightly (0.3331 against 0.3295
on the 44 features; 0.3374 on the association features) and leaves the first pick
unchanged, as it should. Adding which words were already picked does better still
($R^2$ 0.357). But it needs about 130 model evaluations per clue, too slow to play.

In games, the pick-order listener with the win-probability rule (and a one-move
lookahead at the opponent's reply) wins 57.6\% against the incumbent (27 boards
against 13, $p = 0.038$). Compared with the same spymaster without the pick-order
features, though, the difference is $+2.4$ points with a 95\% interval of
$[-4.7, +9.4]$: nothing detectable.

\section{Summary table}
\label{sec:table}

All rows: the same 100 held-out boards, gpt-oss guessing, against the September
incumbent (\texttt{learned\_listener}). $R^2$ is on the Sonnet 5.5 held-out
guesses.

\begin{center}
\small
\begin{tabular}{lcccc}
\toprule
model & $R^2$ & win\% & boards won both ways & sign $p$ \\
\midrule
Gaussian ($\sigma = 1.5$) & -- & \textcolor{red}{TBD} & \textcolor{red}{TBD} & \textcolor{red}{TBD} \\
GBT listener (incumbent) & 0.3205 & 50 & -- & -- \\
\quad + per-pick temperatures & 0.3271 & 44.3 & 8 vs 19 & 0.05 \\
association features & 0.3293 & 53.1 & 23 vs 17 & 0.43 \\
association profile & 0.3309 & 56.7 & 22 vs 9 & 0.029 \\
win probability, association listener & 0.3293 & 59.1 & 28 vs 10 & 0.005 \\
\quad + pick order + reply lookahead & 0.3374 & 57.6 & 27 vs 13 & 0.038 \\
\bottomrule
\end{tabular}
\end{center}

\section{What I found by looking at the clues, and the next step: comfort}

All of the above was optimizing win rate against gpt-oss, and it went up. Then I
played with the best models myself and looked at their clues, and many of them
made no sense to a human. One example: the pick-order model, near the start of a
game, gave \clue{obscure}{4} for \emph{Alien}, \emph{Triangle}, \emph{Spot} and
\emph{Sound}. There is no way a person gets \emph{Triangle} or \emph{Sound} from
that.

I think there are a few causes.
\begin{itemize}
  \item \textbf{Expected value against the modal outcome.} The models maximize
        expected value. People playing Codenames give the clue that does best
        in the outcome they think is most likely. Expected value is probably the
        more correct objective, and it will sometimes give clues where the
        guesser is unsure which word goes with it but that are good on average.
        I don't necessarily want to stop that, even if the spymaster plays
        unlike a human.
  \item \textbf{The win-probability rule chases variance.} After the opponent
        gets 3 words on the first clue, the model thinks it has to catch up with
        a high-variance clue. That usually isn't true this early.
  \item \textbf{The training data never says ``no''.} The listener was trained on
        gpt-oss ranking all 25 words. A ranking always has a 3rd and 4th word,
        so the model learns how gpt-oss orders words it has no real reason to
        pick. A clue for 4 then looks like it buys extra guesses that are close
        to free. Humans don't rank 25 words. They see a clue as pointing at a
        word or not, and they stop when nothing is left. I think we were
        optimizing $R^2$ on later picks that only reflect gpt-oss's habits.
\end{itemize}

So the next step is a measure closer to how the model would do with people. I'm
calling it \emph{comfort}: how clear it is which words a clue is meant for. Win
rate stays; comfort is the second axis.

\paragraph{A different way of asking.} The guesser prompt changes to two separate
calls. The first shows the board and the clue \emph{without} the number. It asks
which words the clue clearly points at (``guess''), which it would still take as
a second or third word if asked for more (``stretch''), and which it wouldn't
guess at all. The answer is a set, in board order, with no ranking. The second
call ranks only the words in that set, now with the number. The guesser guesses
only words in its set and stops when they run out, so it can stop early or pass.

\paragraph{A different way of training the listener.} The same set labels train
a new listener with a STOP option. At each pick it chooses among the words left
and STOP, so it predicts when a guesser would stop instead of assuming it always
uses the full number. On top of it, an \emph{overpromise} penalty charges each
promised word the guesser doesn't attempt (0.5 or 1 word each), and the number is
chosen by the objective.

\paragraph{Measuring comfort.} For every clue I take the set of words gpt-oss
says it clearly points at, which the first call already provides, and compute
$\min(\text{set size}, N)/N$. A score of 1 means the clear words always covered
the number. As a check I asked gpt-oss separately which words it was \emph{sure}
the clue was meant for, on 100 clues per model. 93\% of those words were the
team's, and the two measures rank the models almost identically (Spearman 0.97).

\paragraph{First results.} I reran the menu of current models with the two-step
guesser, against the association-profile incumbent:

\begin{center}
\small
\begin{tabular}{lccc}
\toprule
model & win\% & comfort & mean number \\
\midrule
win probability, association listener & 54.5 & 0.756 & 2.62 \\
overpromise 0.5, strict STOP listener & 52.0 & 0.929 & 1.84 \\
association profile (incumbent) & 51.2$^*$ & 0.839 & 2.20 \\
overpromise 1.0, strict STOP listener & 45.5 & 0.944 & 1.73 \\
STOP listener, loose & 49.0 & 0.781 & 2.50 \\
\bottomrule
\end{tabular}
\end{center}
{\small $^*$ The incumbent's average across all its pairings.}

No model beats the incumbent significantly on wins under the new guesser. The
comfort differences, however, are large and clear. The overpromise model at 0.5
is more comfortable than the incumbent ($t = -8.2$ on the per-board early-stop
rate) and gives up nothing measurable in wins. The STOP listener trained on the
looser set lost to the assassin 19 times against 4 ($p = 0.002$): it expected
the guesser to keep going onto weak words, and this guesser doesn't.

\paragraph{Next.}
\begin{itemize}
  \item Combine the two ends of the trade-off: the win-probability rule with the
        overpromise penalty, converted into win probability at the current
        score. This is running now.
  \item A ``clarity knob'': among clues within $\varepsilon$ of the best win
        probability, pick the clearest, and sweep $\varepsilon$.
  \item Check the comfortable models against real people, which is the point of
        all of this.
\end{itemize}

\end{document}
